% Part: sets-functions-relations
% Chapter: arithmetization
% Section: cuts
%
\documentclass[../../../include/open-logic-section]{subfiles}

\begin{document}

\olfileid{sfr}{arith}{cuts}
\olsection{$\Rat$ توں $\Real$ ول}

اصل وچ تکمیلیت دی خاصیت ایہ وکھاؤندی اے کہ حقیقی خط دا ہر نقطہ
$\alpha$ اوس خط نوں بالکل دو حصیاں وچ ونڈدا اے: اک اوہ جیہدی سب توں
چھوٹی بالائی حد $\alpha$ اے، تے دوجا اوہ جیہدی سب توں وڈی زیریں حد
$\alpha$ اے۔ ناطق عدداں توں حقیقی عدد \emph{بنان} لئی، ڈیڈیکائنڈ نے
سجھایا کہ حقیقی عدداں نوں صرف اوہ \emph{کٹ} سمجھ لیا جاوے جیہڑے ناطق
عدداں نوں ونڈدے نیں۔ یعنی اسی $\sqrt{2}$ نوں اوس \emph{کٹ} نال
شناخت کردے آں جیہڑا $< \sqrt{2}$ ناطق عدداں نوں $> \sqrt{2}$ ناطق
عدداں توں وکھ کردا اے۔

آؤ ایس گل نوں صاف کر لئیے۔ جے اسی ناطق عدداں نوں دو حصیاں وچ ونڈئیے،
تاں صرف اوہدے \emph{نچلے} حصے اُتے غور کر کے اسی اپنی کیتی ونڈ دی
اکّی شناخت کر سکدے آں۔ سو پوری درستی نال اسی ایہ تعریف دیندے آں:

\begin{defn}[کٹ]
اک \emph{کٹ} $\alpha$ ناطق عدداں دا کوئی وی غیر خالی، حقیقی مُڈھلا
ٹکڑا اے جیہدا کوئی سب توں وڈا رکن نہیں۔ یعنی $\alpha$ اک کٹ اُدوں تے
صرف اُدوں اے جدوں:
\begin{enumerate}
	\item \emph{غیر خالی، حقیقی}: $\emptyset \neq \alpha \subsetneq \Rat$
	\item \emph{مُڈھلا}: ہر $p,q \in \Rat$ لئی: جے $p < q \in \alpha$ تاں $p \in \alpha$
	\item \emph{کوئی سب توں وڈا رکن نہیں}: ہر $p \in \alpha$ لئی کوئی $q \in \alpha$ اے جیہدے لئی $p < q$
\end{enumerate}
فیر $\Real$ کٹاں دا سیٹ اے۔
\end{defn}

سو ہُن اسی آکھ سکدے آں کہ $\sqrt{2} = \Setabs{p \in \Rat}{p^2 < 2\text{
یا }p < 0}$۔ بے شک ساڈے نوں جانچنا پینا اے کہ ایہ واقعی اک کٹ
\emph{اے}، پر اسی ایہ کم \olref[check]{sec} ول بھیج دے آں۔

پہلے وانگ، کجھ ہستیاں نوں تعریف کرن توں بعد اگلا کم اوہناں اُتے بنیادی
فنکشن تے تعلق تعریف کرنا اے۔ اسی اک سوکھے تعلق توں شروع کردے آں:
\begin{align*}
	\alpha \leq \beta \text{ اُدوں تے صرف اُدوں جدوں }\alpha \subseteq \beta
\end{align*}
ترتیب دی ایہ تعریف سانوں مرکزی نتیجہ \emph{بیان} کرن دی سہولت دیندی اے:
کٹاں دے سیٹ کول تکمیلیت دی خاصیت اے۔ پوری طرح کھول کے بیان ایہ شکل
رکھدا اے۔ جے $S$ کٹاں دا اک غیر خالی سیٹ اے جیہدی کوئی بالائی حد اے،
تاں $S$ دی اک سب توں چھوٹی بالائی حد اے۔ ہور تفصیل نال: اک کٹ
$\lambda$ اے جیہڑا $S$ دی بالائی حد اے، یعنی $(\forall \alpha \in S)\alpha \subseteq
\lambda$، تے $\lambda$ ایسا سب توں چھوٹا کٹ اے، یعنی $(\forall \beta \in \Real)((\forall \alpha \in S)\alpha \subseteq \beta \lif \lambda \subseteq \beta)$۔ ہُن
نتیجے دا ثبوت ایہ اے:

\begin{thm}\ollabel{realcompleteness}
کٹاں دے سیٹ کول تکمیلیت دی خاصیت اے۔
\end{thm}

\begin{proof}
$S$ کٹاں دا کوئی وی غیر خالی سیٹ ہووے جیہدی بالائی حد اے۔ مقرر کرو کہ
$\lambda = \bigcup S$۔
%\Setabs{p \in \Rat}{(\exists \alpha \in S)p \in \alpha}$$
پہلاں اسی دعویٰ کردے آں کہ $\lambda$ اک کٹ اے:
\begin{enumerate}
\item کیوں جو $S$ غیر خالی اے، ایس لئی گھٹ توں گھٹ اک کٹ $S$ وچ اے، سو
$\emptyset \neq \lambda$۔
% PNBARI-006 ماخذ دی درستی: ایتھے کٹ دی موجودگی بالائی حد توں نہیں بلکہ S دے غیر خالی ہون توں نکلدی اے۔
کیوں جو $S$ کٹاں دا سیٹ اے، $\lambda \subseteq \Rat$۔ کیوں جو $S$ دی
بالائی حد اے، کوئی $p \in \Rat$ ہر کٹ $\alpha \in S$ وچوں غائب اے۔
سو $p\notin \lambda$، تے ایس لئی $\lambda \subsetneq \Rat$۔
\item فرض کرو $p < q \in \lambda$۔ سو کوئی $\alpha \in S$ اے جیہدے
لئی $q \in \alpha$۔ کیوں جو $\alpha$ اک کٹ اے، $p \in \alpha$۔ سو
$p \in \lambda$۔
\item فرض کرو $p \in \lambda$۔ سو کوئی $\alpha \in S$ اے جیہدے لئی
$p \in \alpha$۔ کیوں جو $\alpha$ اک کٹ اے، کوئی $q \in \alpha$ اے
جیہدے لئی $p < q$۔ سو $q \in \lambda$۔
\end{enumerate}
ایہ دعویٰ ثابت کردا اے۔ نال ای صاف اے کہ $(\forall \alpha \in S)\alpha
\subseteq \bigcup S = \lambda$، یعنی $\lambda$، $S$ دی بالائی حد اے۔
سو ہُن فرض کرو کہ $\beta \in \mathbb{R}$ وی اک بالائی حد اے، یعنی
$(\forall \alpha \in S)\alpha \subseteq \beta$۔ ہر $p \in \Rat$ لئی،
جے $p \in \lambda$، تاں کوئی $\alpha \in S$ اے جیہدے لئی $p \in
\alpha$، سو $p \in \beta$۔ عام کر کے $\lambda \subseteq \beta$۔ سو
$\lambda$، $S$ دی \emph{سب توں چھوٹی} بالائی حد اے۔
\end{proof}

سو ساڈے کول ہُن ہستیاں دا اک جتھا اے جیہڑا تکمیلیت دی خاصیت پوری کردا
اے۔ ایس گل نوں بیان کرن دا اک طریقہ ایہ اے: ساڈے کٹاں وچ کوئی ”خلا“
نہیں۔ (سو ناطق عدداں دی بجائے حقیقی عدداں دے ہور ”کٹ“ لین نال کوئی
دلچسپ نویں شے نہیں ملے گی۔)

اگے ساڈے نوں حقیقی عدداں اُتے کجھ عمل تعریف کرنے پینے نیں۔ اسی ایہ مقرر
کر کے ناطق عدداں نوں حقیقی عدداں وچ سماؤندے آں کہ $p_\Real =
\Setabs{q \in \Rat}{q < p}$، ہر $p \in \Rat$ لئی۔ فیر اسی تعریف کردے آں:
\begin{align*}
	\alpha + \beta &= \Setabs{p + q}{p \in \alpha \land q \in \beta}\\
%	\alpha - \beta &\defis \Setabs{p - q}{p \in \alpha \land q \in \Rat \setminus \beta}\\
	\alpha \times \beta &=
	\Setabs{p \times q}{0 \leq p \in \alpha \land 0 \leq q \in \beta} \cup 0_\Real & \text{جے }\alpha, \beta \geq 0_\Real
%	\alpha \div \beta &\defis
%	\Setabs{\nicefrac{p}{q}}{p \in \alpha \land q \in \Rat \setminus \beta}  & \text{جے }\alpha \geq 0_\Real, \beta > 0_\Real
\end{align*}
% PNBARI-007 ماخذ دی درستی: ناطق p دی جاگذاری p_Real اے؛ ایس لئی ضرب وچ حقیقی صفر وی 0_Real اے، اک واری آیا 0^R نہیں۔
ضرب دے دوجے معاملے سنبھالن لئی، پہلاں مقرر کرو: %کہ $0_\Real \times \alpha = 0_\Real = \alpha \times 0_\Real$، تے شامل کرو:
\begin{align*}
	-\alpha &= \Setabs{p - q}{p < 0 \land q \notin \alpha}
\end{align*}
تے فیر ایہ مقرر کرو:
\begin{align*}
	\alpha \times \beta &\defis
	\begin{cases}
		\mathord{-}\alpha \times \mathord{-}\beta &\text{جے }\alpha < 0_\Real\text{ تے }\beta < 0_\Real\\
		\mathord{-}(\mathord{-}\alpha \times \beta) &\text{جے }\alpha < 0_\Real \text{ تے }\beta > 0_\Real\\
		\mathord{-}(\alpha \times \mathord{-}\beta) &\text{جے }\alpha > 0_\Real \text{ تے }\beta < 0_\Real
	\end{cases}
\end{align*}
فیر ساڈے نوں جانچنا پینا اے کہ ایہناں وچوں ہر تعریف ہمیشہ اک کٹ دیندی
اے۔ آخر وچ ساڈے نوں اک سوکھا (پر لمبا چوڑا) مظاہرہ کرنا پینا اے کہ
ایویں تعریف کیتے کٹ عین اوہو جیہا ورتاؤ کردے نیں جیہو جیہا اوہناں نوں
کرنا چاہیدا اے۔ پر اسی ایہ سارا کم \olref[check]{sec} ول بھیج دے آں۔
\end{document}
